\documentclass[xcolor=table,dvipsnames,table]{beamer}
\mode<presentation>
\usetheme{boxes}
\setbeamertemplate{navigation symbols}{}
% http://www.latex-community.org/forum/viewtopic.php?f=4&t=6694
\setbeamertemplate{navigation symbols}{\raisebox{5pt}{\makebox[\paperwidth]{\hfill\makebox[10pt]{\scriptsize\insertframenumber\vspace{1ex}}}}}
%\setbeamertemplate{footline}[frame number]
\setbeamertemplate{blocks}[shadow=false]
%\setbeamercolor*{block title}{fg=structure,bg=RoyalBlue!10}
\setbeamercolor*{block title example}{fg=structure,bg=RoyalBlue!10}
%\setbeamercolor*{block title example}{fg=BrickRed,bg=Goldenrod!10}
\setbeamercolor*{block title alerted}{fg=white,bg=black}
\addtobeamertemplate{block begin}{\pgfsetfillopacity{0.8}}{\pgfsetfillopacity{1}}
%\rowcolors{0}{RoyalBlue!20}{RoyalBlue!5}

%\DeclareGraphicsRule{*}{mps}{*}{}

\usepackage{latexsym}
\usepackage{hyperref}
\usepackage{tikz}
\usetikzlibrary{calc,shapes,arrows,shadows,shapes.callouts,shapes.arrows,chains,positioning,trees}
\usepackage{solution}
\usepackage{calc}
\usepackage{pifont}
\usepackage{algorithmic}

\newcommand{\cmark}{\ding{51}}
\newcommand{\xmark}{\ding{55}}

\newcounter{mycallout}

\newcommand{\callouts}[3]{%
  \stepcounter{mycallout}
  \tikz[remember picture,baseline]{\node[anchor=base,inner sep=0,outer sep=0]%
    (\themycallout) {\colorbox{#1!20}{#3}};\pause\node[overlay,rectangle callout,%
    callout relative pointer={(0cm,0.5cm)},fill=#1!20] at ($(\themycallout.south)+(-0cm,-0.7cm)$){#2};}%
    }%

\raggedright

\newcount\lecturecount
\lecturecount=0
\AtBeginLecture{%
    \advance\lecturecount by 1
    \date{}
    \begin{frame}
    \begin{center}
    \titlepage
    \ifnum\lecturecount=1
    Part \the\lecturecount: \insertlecture
    \else
    Part \the\lecturecount: \insertlecture
    \fi
    \end{center}
    \end{frame}
}

\addtobeamertemplate{block begin}{\setlength\abovedisplayskip{0pt}}

%\newcommand{\example}[1]{{\color{BrickRed!50}{#1}}}
\newcommand{\maths}[1]{{\color{RoyalBlue!50}{#1}}}
\newcommand{\reference}[1]{{\color{RoyalBlue!30}\tiny [from #1]}}
\newcommand{\koehnref}{\reference{\href{http://www.statmt.org/book}{P.Koehn SMT book slides}}}


\ifx\pdfoutput\undefined
  \usepackage{graphicx}
\else
  %\usepackage[pdftex]{graphicx}
  \DeclareGraphicsRule{*}{mps}{*}{}
\fi
\begin{document}

\title{\color{RoyalBlue!20}Natural Language Processing}

\author{Anoop Sarkar \\ {\tt anoopsarkar.github.io/nlp-class}}
\institute{Simon Fraser University}
%\date{}
     
{
\setbeamertemplate{navigation symbols}{}
\addtocounter{framenumber}{-1}
\begin{frame}
\begin{center}
\vspace{8mm}
\includegraphics[scale=0.35]{figures/natlang-cky-logo}
\end{center}
\titlepage
\end{frame}
}



\lecture{Ambiguity}{}
\section{Context Free Grammars and Ambiguity}

\begin{frame}[fragile]
\frametitle{Context Free Grammars and Ambiguity}
\[
\begin{array}{ccc}
 S & \rightarrow & NP~VP  \\
 VP & \rightarrow & V~NP   \\
 VP & \rightarrow & VP~PP \\
 PP & \rightarrow & P~NP   \\
 NP & \rightarrow & NP~PP  \\
 NP & \rightarrow & Calvin  \\
 NP & \rightarrow & monsters  \\
 NP & \rightarrow & school  \\
 V & \rightarrow & imagined \\
 P & \rightarrow & in     
\end{array}
\]
What is the analysis using the above grammar for:\\
{\em Calvin imagined monsters in school}
\end{frame}

\begin{frame}[fragile]
\frametitle{Context Free Grammars and Ambiguity}
{\em Calvin imagined monsters in school}
\begin{verbatim}
(S (NP Calvin)
   (VP (V imagined)
       (NP (NP monsters)
           (PP (P in)
               (NP school)))))

(S (NP Calvin)
   (VP (VP (V imagined)
           (NP monsters))
       (PP (P in)
           (NP school))))
\end{verbatim}
Which one is more plausible?
\end{frame}

\begin{frame}[fragile]
\frametitle{Ambiguity Kills (your parser)}
{\em natural language learning course} \\
(run demos/parsing-ambiguity.py)
\begin{verbatim}
((natural language) (learning course))
(((natural language) learning) course)
((natural (language learning)) course)
(natural (language (learning course)))
(natural ((language learning) course))
\end{verbatim}
\begin{itemize}
\item Some difficult issues:
\begin{itemize}
\item Which one is more plausible?
\item How many analyses for a given input? 
\item Computational complexity of parsing language
\end{itemize}
\end{itemize}
\end{frame}

\begin{frame}
\frametitle{Treebanks}
\begin{itemize}
\item What is the CFG that can be extracted from this single tree: 
\begin{tabbing}
0123\=4567\=8901\=2345\=6789\=0123\=4567 \kill
(S \>(NP (Det the) (NP man)) \\
\> (VP \> (VP \> (V played) \\
\> \> \> (NP (Det a) (NP game))) \\
\> \> (PP \>(P with) \\
\> \> \> (NP (Det the) (NP dog)))))
\end{tabbing}
\end{itemize}
\end{frame}

\begin{frame}
\frametitle{PCFG}
{\small
\[
\begin{array}{cccc}
 S & \rightarrow & NP~VP  & c=1  \\
 NP & \rightarrow & Det~NP & c=3  \\
 NP & \rightarrow & man  & c=1   \\
 NP & \rightarrow & game  & c=1   \\
 NP & \rightarrow & dog  & c=1   \\
 VP & \rightarrow & VP~PP  & c=1  \\
 VP & \rightarrow & V~NP  & c=1  \\
 PP & \rightarrow & P~NP  & c=1  \\
 Det & \rightarrow & the  & c=2  \\
 Det & \rightarrow & a  & c=1   \\
 V & \rightarrow & played & c=1  \\
 P & \rightarrow & with & c=1 
\end{array}
\] 
}
\begin{itemize}
\item We can do this with multiple trees. Simply count occurrences of CFG rules over all the trees. 
\item A repository of such trees labelled by a human is called a TreeBank.
\end{itemize}
\end{frame}

\begin{frame}
\frametitle{Ambiguity}
\begin{itemize}
  \item Part of Speech ambiguity\\
        {\tt saw} $\rightarrow$ {\color{red}noun}\\
        {\tt saw} $\rightarrow$ {\color{red}verb}
  \item Structural ambiguity: Prepositional Phrases\\
        {\tt I saw (the man) with the telescope}\\
        {\tt I saw (the man with the telescope)}
  \item Structural ambiguity: Coordination\\
        {\tt a program to promote safety in ((trucks) and (minivans))}\\
        {\tt a program to promote ((safety in trucks) and (minivans))}\\
        {\tt ((a program to promote safety in trucks) and (minivans))}
\end{itemize}

\end{frame}

\begin{frame}[fragile]
\frametitle{Ambiguity $\leftarrow$ attachment choice in alternative parses}
\includegraphics[scale=.6]{figures/figures.18}
\end{frame}

\begin{frame}
\frametitle{Ambiguity in Prepositional Phrases}
\begin{itemize}
\item<1-> noun attach: {\em I bought the shirt with pockets}
\item<2-> verb attach: {\em I washed the shirt with soap}
\item<3-> As in the case of other attachment decisions in parsing: it depends on the meaning of the entire sentence -- needs world knowledge, etc.
\item<4-> Maybe there is a simpler solution: we can attempt to solve it using heuristics or associations between words
\end{itemize}

\end{frame}

\begin{frame}
\frametitle{Structure Based Ambiguity Resolution}
\begin{itemize}
  \item Right association: a constituent (NP or PP) tends to attach to
  another constituent immediately to its right (Kimball 1973)
  \item Minimal attachment: a constituent tends to attach to an
  existing non-terminal using the fewest additional syntactic nodes
  (Frazier 1978)
  \item These two principles make opposite predictions for
  prepositional phrase attachment
  \item Consider the grammar:
\begin{eqnarray}
VP &\rightarrow& V\ NP\ PP \label{rule:VP}\\
NP &\rightarrow& NP\ PP \label{rule:NP}
\end{eqnarray}  
  for input: {\em I [$_{VP}$ saw [$_{NP}$ the man $\ldots$ [$_{PP}$ with
  the telescope ]}, \\
  RA predicts that the PP attaches to the NP, i.e. use rule (\ref{rule:NP}), \\
  and MA predicts V attachment, i.e. use rule (\ref{rule:VP})
\end{itemize}

\end{frame}

\begin{frame}
\frametitle{Structure Based Ambiguity Resolution}
\begin{itemize}
  \item Garden-paths look structural:\\
  {\rm\it The emergency crews hate most is
  domestic violence}
  \item Neither MA or RA account for more than 55\% of the cases in
  real text 
  \item Psycholinguistic experiments using eyetracking show that humans resolve
  ambiguities as soon as possible in the left to right sequence using
  the words to disambiguate
  \item Garden-paths are caused by a combination of lexical and structural effects:\\ 
  {\rm\it The flowers delivered for the patient arrived}
\end{itemize}

\end{frame}

\begin{frame}
\frametitle{Ambiguity Resolution: Prepositional Phrases in English}
  \begin{itemize}
  \item Learning Prepositional Phrase Attachment: Annotated Data
\begin{tabular}{|cccc|c|}
\hline
v    &     n1   &      p & n2 &       Attachment\\
\hline
join   &   board    &   as &  director & V \\
is       & chairman  &  of  & N.V.    & N \\
using  &   crocidolite & in  & filters & V \\
bring  &   attention  & to  & problem &  V \\ 
is      &  asbestos   & in &  products & N \\
making &   paper    &   for & filters & N \\
including & three     &  with & cancer &  N \\
$\vdots$ & $\vdots$ & $\vdots$ & $\vdots$ & $\vdots$ \\
\hline
\end{tabular}
  \end{itemize}

\end{frame}

\begin{frame}
\frametitle{Prepositional Phrase Attachment}
\begin{tabular}{|l|c|}  \hline
Method & Accuracy \\ \hline
Always noun attachment & 59.0 \\
Most likely for each preposition & 72.2 \\
Average Human (4 head words only) & 88.2 \\
Average Human (whole sentence) & 93.2 \\  \hline
\end{tabular}
\end{frame}

\begin{frame}
\frametitle{Some other studies}
  \begin{itemize}
%    \item {\bf Brill and Resnik, 1994}: \\
%use transformation based learning for PP attachment\\
%80.8\% with words; with Wordnet classes: 81.8\%
  \item {\bf Toutanova, Manning, and Ng, 2004}: 87.54\% using some external knowledge (word classes)
  \item {\bf Merlo, Crocker and Berthouzoz, 1997}: test on multiple PPs
  \begin{itemize}
	\item generalize disambiguation of 1 PP to 2-3 PPs
	\item 14 structures possible for 3PPs assuming a single verb
	\item all 14 are attested in the Penn WSJ Treebank
	\item 1PP: 84.3\% \ \ 2PP: 69.6\% \ \ 3PP: 43.6\% 
	\item {\bf This experiment is still only part of the real problem faced in parsing English}
	\item Other sources of ambiguity in other languages
  \end{itemize}
  \end{itemize}

\end{frame}

\end{document}
